\subsection{A validity check on recorded sale prices}
\label{sec:saleprice}

The outcome throughout is the listing's asking price, which an agent sets together
with the description and which sits well above the price a property eventually
clears, so a natural question is whether the text effect we estimate is a property
of the market's valuation or an artifact of the agent pricing their own language.
Four of the twelve metropolitan areas, Philadelphia, Chicago, Washington, and New
York, publish recorded transaction prices through their county assessor or
recorder, and for these we recover the realized sale price and its date, match
each listing to the recorded sale of the parcel it occupies, and re-estimate the
text effect on the transaction rather than the ask. The treatment and the
structured and spatial controls are held fixed, and we add sale-quarter fixed
effects because recorded sales span several years in nominal dollars. The asking
price in these markets runs between a fifth and ninety percent above the eventual
sale, so the two outcomes are far from interchangeable.

The effect survives the change of outcome. Estimated on the realized sale price it
is positive and individually significant in all four markets, at $+0.358$ in
Philadelphia, $+0.107$ in Chicago, $+0.070$ in Washington, and $+0.042$ in New
York, against listing-price estimates on the same matched sample of $+0.472$,
$+0.431$, $+0.151$, and $+0.039$ reported alongside in Table~\ref{tab:saleprice}.
The sale-price effect is smaller than the listing-price effect in three of the
four markets, by about a quarter in Philadelphia and by more in Chicago, which is
the signature of the asking-price simultaneity we are explicit about, since part
of the association between language and the ask is the agent's own pricing rather
than the market's. New York, whose effect is near zero on either outcome, is
unchanged and remains the panel's exception. We therefore read the listing-price
estimates as an upper bound on the price effect of a property's language, with the
realized-transaction effect positive but materially smaller, and we report the
per-market replication rather than the four-market pooled mean, which a
Hartung-Knapp interval leaves wide given the heterogeneity across so few markets.
That the effect is present at all on the transaction the market actually cleared,
in every market where that transaction is public, is the strongest evidence the
available data affords that the listing-price association is not merely an artifact
of how agents price their own prose.

\begin{table}[H]
\centering
\caption{Sale-price validity check. Per-market DML effect of the leading text
direction on log listing price and on log recorded sale price, estimated on the
matched sample with identical treatment and controls; the sale-price column adds
sale-quarter fixed effects. The final column is the median ratio of asking to
recorded sale price. Recorded sales are from the Philadelphia OPA, the Cook County
Assessor, the DC CAMA system, and the NYC Department of Finance.}
\label{tab:saleprice}
\small
\begin{tabular}{lrrrr}
\toprule
Market & $n$ & $\hat\theta$ list price & $\hat\theta$ sale price & list/sale \\
\midrule
Philadelphia & $6{,}604$ & $+0.472\ (0.012)$ & $+0.358\ (0.014)$ & $1.90$ \\
Chicago      & $4{,}746$ & $+0.431\ (0.018)$ & $+0.107\ (0.017)$ & $1.53$ \\
Washington   & $2{,}114$ & $+0.151\ (0.017)$ & $+0.070\ (0.012)$ & $1.22$ \\
New York     & $6{,}436$ & $+0.039\ (0.015)$ & $+0.042\ (0.015)$ & $1.24$ \\
\bottomrule
\end{tabular}
\end{table}
